\begin{center}
\textsc{Gauss} to \textsc{Gerling}.\quad Göttingen, 19~July 1827.
\end{center}

\par\noindent\hspace*{\parindent}\mbox{.\ .\ .\ .\ .\ .}\quad Although the weather was less favourable to my astronomical work
than I had wished, I nevertheless believe I have determined the amplitude
of the arc between the observatories at Göttingen
and Altona to within a very small fraction of an arcsecond. I used
a very large number of stars and
made some 900 observations in all. There remains a difference of
\par\clearpage
\noindent\null\hfill{\small\textbf{ON THE HANOVER TRIANGULATION.}}\hfill\makebox[0pt][r]{\fontsize{13}{16}\selectfont 391}\par\medskip
\noindent \(5^{\prime\prime}\) from the calculation based on the triangles using \textsc{Walbeck}'s dimensions of the Earth,
and so it would have been very interesting to have another point
in between. With the sector it was impossible under the prevailing circumstances.
But perhaps we shall observe together in Celle at some future time, with a transit
instrument to be set up in the prime vertical, with \textsc{Schumacher} observing
beforehand in Altona and me afterwards in \mbox{Göttingen.\quad\mbox{.\ .\ .\ .\ .\ .}}\quad Before my
journey I was also sent all the Bavarian triangles,
which are linked to mine by \mbox{yours.\quad\mbox{.\ .\ .\ .\ .\ .}}
\par
